% R1: validation (Vaishak). Numbers from scripts/aggregate_seeds.py over 5 seeds (one CCR
% run); tables from scripts/make_methods_tables.py. Needs: booktabs, graphicx.

\section{Validation Protocol}\label{sec:validation}

Windows that overlap by half and come from the same 15-minute recording are almost copies of
each other. A random split over windows therefore puts near-duplicates of every test window in
the training set. We compare three splits:
\begin{itemize}
  \item \textbf{Naive}: six stratified folds over windows. It leaks, and is reported only to
        measure that leakage.
  \item \textbf{Grouped}: each class's six recordings are dealt one per fold, so every test fold
        holds one whole, unseen recording per class. This is our main protocol.
  \item \textbf{Leave one condition out}: fold $k$ holds out every class's recording at
        operating condition $k$.
\end{itemize}
Because each class is recorded once per condition, every grouped test recording is also from a
condition that the training fold never saw for that class. Scores are macro-F1 averaged over the
six folds; chance is $1/6 \approx 0.17$.

\textbf{Seeds.} Fold-to-fold spread is large (up to $\pm$0.2 macro-F1), so a single fold
assignment can make one model look best by luck. Every experiment is repeated over five seeds.
A seed changes the grouped fold assignment, the random-forest seed, the Kernel PCA subsample and
the corruption noise. We report the mean and standard deviation over seeds, all from one run on
one machine.

\textbf{Separability on held-out recordings.} A grouped test fold has a single recording per
class, so the usual Fisher ratio or silhouette computed on it measures how distinct six
individual recordings are, not how well the classes generalise. We therefore score held-out
separability with a transfer Fisher ratio: the between-class scatter comes from the training
class means, and the within-class scatter is the spread of held-out windows around those same
training means, Ledoit--Wolf shrunk. Fisher $J$ is compared only between spaces of equal
dimension.

\begin{table}[t]
  \centering
  \caption{Best model under each split (macro-F1, mean $\pm$ std over 5 seeds). The
  leave-one-condition-out folds do not depend on the seed.}
  \label{tab:splits}
  \begin{tabular}{llc}
\toprule
Split & Best model & macro-F1 \\
\midrule
Naive (random windows) & 25 of 30 models within 0.005 & 1.00 $\pm$ 0.00 \\
Leave one condition out & Raw 126-D + Lin.\ SVM & 0.56 $\pm$ 0.00 \\
Grouped by recording & Raw 126-D + Lin.\ SVM & 0.47 $\pm$ 0.03 \\
\bottomrule
\end{tabular}

\end{table}

\begin{table*}[t]
  \centering
  \caption{Grouped cross-validation: macro-F1 of every front-end and classifier (mean
  $\pm$ std over 5 seeds; chance 0.17). Bold: best.}
  \label{tab:grid}
  \begin{tabular}{lcccccc}
\toprule
Front-end & $k$-NN & $k$-NN (Mah.) & Lin.\ SVM & RBF SVM & RF & GNB \\
\midrule
Raw 126-D & 0.28 $\pm$ 0.03 & 0.25 $\pm$ 0.04 & \textbf{0.47 $\pm$ 0.03} & 0.35 $\pm$ 0.04 & 0.39 $\pm$ 0.05 & 0.27 $\pm$ 0.02 \\
PCA (95\% var.) & 0.29 $\pm$ 0.01 & 0.30 $\pm$ 0.03 & 0.41 $\pm$ 0.02 & 0.38 $\pm$ 0.03 & 0.40 $\pm$ 0.03 & 0.27 $\pm$ 0.06 \\
PCA-5 & 0.28 $\pm$ 0.03 & 0.23 $\pm$ 0.01 & 0.31 $\pm$ 0.04 & 0.36 $\pm$ 0.04 & 0.29 $\pm$ 0.04 & 0.24 $\pm$ 0.03 \\
LDA-5 & 0.45 $\pm$ 0.06 & 0.42 $\pm$ 0.06 & 0.44 $\pm$ 0.07 & 0.43 $\pm$ 0.05 & 0.42 $\pm$ 0.05 & 0.43 $\pm$ 0.06 \\
Kernel PCA-5 & 0.15 $\pm$ 0.01 & 0.16 $\pm$ 0.02 & 0.28 $\pm$ 0.05 & 0.20 $\pm$ 0.02 & 0.21 $\pm$ 0.06 & 0.21 $\pm$ 0.03 \\
\bottomrule
\end{tabular}

\end{table*}

\textbf{Results.} The leakage is the first result (Table~\ref{tab:splits}). Every model scores
about 1.00 under the naive split and at most 0.47 when whole recordings are held out. Leaving one
condition out scores higher than the grouped split (0.56), probably because all six classes
are then tested on the same unseen condition, which preserves their differences relative to
each other.

Across the grid (Table~\ref{tab:grid}) the best model is the linear SVM on the raw features
($0.47 \pm 0.03$). LDA-5 is uniformly decent with every classifier (0.42--0.45) but is not ahead
once seeds are averaged. With a single seed, LDA-5 with a random forest reached 0.49, which is
the kind of luck the seed average removes. LDA's in-sample Fisher $J$ at $k=5$ is about 490,
while its held-out $J$ is $15.1 \pm 4.7$, close to PCA-5 ($10.1 \pm 3.5$): most of what LDA
separates on its training fold is which recording a window came from. Kernel PCA-5 is the
weakest front-end ($0.28 \pm 0.05$ at best).

Per class, B1 is recognised almost perfectly (recall $0.95 \pm 0.03$ for the best model),
while B3 (0.31), R (0.41) and B2 (0.43) remain hard. Section~\ref{sec:separability} explains why.
